\section{Information}\label{sec:info}

\paragraph{Frozen readers exaggerate the loss of truth information.} A logistic truth probe fitted on original
activations and applied unchanged to SAE reconstructions loses 0.28--0.29 AUROC on Geometry of Truth
(Figure~\ref{fig:m2}). Isotropic noise with the norm of the SAE error produces almost the same drop (0.25--0.26), and a
simpler mass-mean reader loses only 0.02--0.03. Most of the apparent collapse therefore reflects the sensitivity of
that reader to any perturbation of this size rather than information removed by the SAE.

\paragraph{Re-trained readers show a smaller but real loss.} Readers re-trained on the reconstructions are 0.07--0.12
AUROC worse than readers trained on original activations, and no cell is classified as preserved. The label can also be
decoded from the reconstruction error alone (AUROC 0.62--0.68), so part of the truth information remains in the
activation but outside the reconstruction.

\paragraph{Persona information is lost, not merely misread.} On the equal-length RepE-LM prompts the frozen persona
gap is 1.9--3.4 times the gap caused by norm-matched noise, and re-trained readers recover little of it: the
re-trained gap is 0.08--0.22, and its statement-bootstrap interval excludes zero for all 99 base SAEs at layers 1--11.
Part of the loss is generic, since a neutral one-word swap also loses information, but at layers 2--8 the persona gap
exceeds the neutral control by 0.06--0.12. The persona contrast is also distorted geometrically: for each pair, the
reconstructed honest-minus-dishonest difference has a cosine of only 0.35--0.58 with the original difference, even
though the mean difference direction over all pairs is kept better (0.70--0.86). The SAE reconstructs each prompt
reasonably well but their difference poorly, because the two reconstruction errors do not cancel.

\paragraph{The SAE-versus-PCA verdict depends on the constraint.} Matched to the SAE's OWT reconstruction error, PCA
keeps at least as much task information in every cell, and the SAE is significantly worse in 52 of 81 truth cells, 59
of 99 persona cells and 81 of 90 neutral cells. In the pre-registered capacity-matched replication on the 20 trained
family$\times$layer cells (Table~\ref{tab:pcacm}), the SAE is worse in 13 and never better. That comparison, however,
gives PCA 298--519 dense components, whereas each SAE activates about 128 latents per token. Restricted to $k=128$
components, no trained SAE is classified worse and three of 20 cells are better; with $k$ matched to each SAE's
task-level activity ($k=54$--128), eight are better and none worse. Under an approximate post-hoc Holm correction over
the 40 capacity-matched cells, one of the three and four of the eight remain. ``Not worse'' is not equivalence: the
pre-registered $\pm0.02$ equivalence class is met in none of the 40 cells, because the paired intervals are wider than
the margin. Whether an SAE preserves more information than PCA therefore depends on whether reconstruction error or the
per-token active budget is equalised. The persona and neutral tasks were compared at equal error only.

\begin{figure}[t]
  \centering
  \includegraphics[width=\linewidth]{fig_m2_information}
  \caption{Information preservation by layer for the base GPT-2 SAEs (\evlabel{validation}). Dashed: frozen logistic
  AUROC gap; solid: re-trained logistic gap; dotted grey: frozen gap under norm-matched noise; black: PCA matched to the
  ReLU SAE's reconstruction error. Per-architecture values are in Appendix Table~\ref{tab:m2}.}
  \label{fig:m2}
\end{figure}

% generated by scripts/shared/pca_cm_posthoc_multiplicity.py from results/shared/pca_capacity_matched/
\begin{table}[t]
\centering
\small
\caption{PCA-CM (\evlabel{validation}; GoT cities+neg\_cities VAL, 448 rows): SAE vs PCA on the GoT re-fit logistic truth gap (lower is better) under three matching constraints, over the 20 trained family$\times$layer cells (L8, L11; random decoder excluded). ``Better''/``worse'': every seed's paired 95\,\% grouped-bootstrap CI excludes zero (pre-registered rule, no correction across cells). ``Holm'': cells that also survive an approximate post-hoc Holm correction over the 40 capacity-matched cells. Better at $k=128$ (SAE vs PCA gap): L8 R2 ReLU: 0.012 vs 0.058; L8 res-jb: 0.025 vs 0.058; L11 R2 ReLU: 0.076 vs 0.153.}
\label{tab:pcacm}
\begin{adjustbox}{max width=\linewidth}
\begin{tabular}{lccccc}
\toprule
PCA reference & PCA $k$ & SAE better & of which Holm & inconclusive & SAE worse \\
\midrule
Equal OWT error (existing reference) & 298--519 & 0 & -- & 7 & 13 \\
Fixed $k=128$ (per-token active budget) & 128 & 3 & 1 & 17 & 0 \\
$k$ = SAE's GoT task L0 & 54--128 & 8 & 4 & 12 & 0 \\
\bottomrule
\end{tabular}
\end{adjustbox}
\end{table}

